% Auto-generated related-work appendix.
% Source: paper/related-work-cache/sections/*.md (plan order)
% Generation pipeline: HERALD related-work orchestration (paper-readers, quote-verifiers, critics, section-writers).

\appendix
\section*{Appendix: Extended Related Work}

\section{KV-Cache Compression Methods: The Landscape HERALD Sits Inside}

KV-cache compression is now a crowded design space. To position HERALD, which forecasts compression-induced catastrophic failures from per-token logit signals, we map the families of compressors that share one critical property: each is a black-box transform of the cache that HERALD's matched-prefix replay harness can swap in, and whose aggressiveness a runtime controller can in principle modulate. We organise the landscape by what the compressor manipulates: which tokens are kept (eviction), which are read at attention time (sparse attention), how each KV is represented (quantization and low-rank), and how non-resident tokens are reintroduced (retrieval and chunked recomputation).

\textbf{Eviction-based compression.} The dominant family decides, online, which KV pairs to drop. H2O retains "heavy hitter" tokens identified by accumulated attention scores, exploiting the empirical observation that over 95\% of attention concentrates on under 5\% of tokens \cite{h2o}. SnapKV instead uses a prompt-end observation window to vote on which positions each head will attend to during decoding, then snaps the cache to those positions \cite{snapkv}. PyramidKV refines uniform-budget eviction by allocating more cache to lower layers and less to upper layers, exploiting the pyramidal information-funnelling pattern, and reports matching full-cache quality at 12\% retention \cite{pyramidkv}. Ada-KV pushes this finer, allocating budget per attention head in proportion to attention concentration with a theoretical L1 bound on eviction loss, and works as a plug-in to SnapKV and PyramidKV \cite{adakv}. CompressKV identifies semantically-aware retrieval heads and applies layer-adaptive allocation guided by per-layer eviction error \cite{compresskv}. G-KV addresses the observation that the last window's attention patterns are not fully consistent with those of earlier windows, an inconsistency that grows as fewer tokens are retained, and combines local and historical attention into a global score \cite{gkv}. R-KV jointly scores attention importance and key-vector redundancy, motivated by the finding that reasoning models emit 8-14x longer outputs with 5-7x more n-gram repetition, exactly the regime HERALD targets \cite{rkv}. EvicPress goes further and jointly tunes eviction and per-token compression configurations per context, motivated by the observation that no single compression strategy is universally optimal \cite{evicpress}. Many further eviction variants populate this corner of the design space \cite{anchorattention,baklava,dbudgetkv,keepkv,keyvaluecompress,kvcompress,lava,lazyeviction,rap,refreshkv,rethinkingkv,statefulkv,fastkvzip,costoptimalgqa}.

\textbf{Sparse-attention and streaming methods.} A second family keeps the full cache resident but only reads a sparse subset at attention time. StreamingLLM pairs a sliding window with preserved initial-token "attention sinks" and processes up to 4M tokens without retraining \cite{streamingllm}; the attention-sink phenomenon it exploits has since been traced to softmax normalisation, with sigmoid attention eliminating the effect \cite{attentionsink}. Quest dynamically loads only the top-k KV pages per query using min/max page statistics, sustaining lossless accuracy on long-context retrieval while accelerating self-attention \cite{quest}. DuoAttention partitions heads into "retrieval heads" needing full cache and "streaming heads" needing only a constant-length cache, enabling 3.3M-token Llama-3-8B decoding on a single A100 \cite{duoattention}. Native Sparse Attention combines blockwise compression, learnable token selection, and a sliding window into a single hardware-aligned mechanism, giving 11.6x speedup at 64K context with perfect needle-in-a-haystack accuracy \cite{nsa}. RocketKV stages SnapKV-style permanent eviction with hybrid sparse attention during decode, reaching 400x compression ratios \cite{rocketkv}. Expected Attention takes a third route, training-free: it analytically estimates future query distributions from the Gaussian structure of hidden states under RoPE, scoring each KV by its expected future contribution, and operates during both prefill and decode \cite{expectedattention}.

\textbf{Quantization, low-rank, and structural compression.} Orthogonal to which tokens are kept, another family compresses how each KV is stored. Quantization-based methods squeeze keys and values into low-precision representations \cite{atom,kivi,kvquant,quarot,turboquant,zipcache,gear}, low-rank and projection-based methods factor the cache or migrate to MLA-style latents \cite{palu,transmla,kvlinc}, and architectural variants tune cache structure for specific regimes \cite{kitty}. These methods can be composed with eviction or sparse-attention without changing the cache's logical interface.

\textbf{Retrieval and reconstruction.} A fourth family explicitly accepts that compression is lossy and reconstructs missing context. KVzip prunes KVs by maximum chunk-attention activation and reconstructs evicted entries via forward-pass recomputation at 5-15\% per-token overhead, reaching 4-8x compression with sub-1\% accuracy loss across SQuAD, GSM8K, and NIAH \cite{kvzip}. KVzap distils KVzip-style scores into a lightweight surrogate for adaptive thresholded eviction during prefill and decode \cite{kvzap}. Semantic chunking is a related axis: ChunkKV groups consecutive tokens, scores chunks, and reuses chunk indices across layers, achieving 4.5x compression at 99\% of GSM8K quality \cite{chunkkv}. Retrieval-augmented and rehydration approaches generalise this idea, swapping cache for an external store \cite{retrievalattention,shadowkv,cacheblend}.

\textbf{What HERALD treats as a black box.} HERALD does not propose a new compressor or modify any existing one. It assumes only that a compressor exposes a `press(model)` context with a configurable compression ratio, after which generation proceeds with `output_scores=True`. Every method above, eviction policies, query-aware sparse readers, quantizers, low-rank factorisations, chunked retrievers, satisfies this interface; conversely, methods that re-tune the cache architecturally at training time \cite{duoattention,nsa,transmla} expose the same runtime surface once deployed. This is what makes the matched-prefix replay measurement methodology of phase 0 portable across families: identical prompts and decoding state, only the compressor swapped. It is also what makes the closed-loop ambition of later phases tractable: any of these compressors can be the substrate whose ratio a runtime controller modulates, with HERALD's per-token hazard score as the control signal that fires before catastrophic failure manifests.

\section{Runtime Quality Signals: Entropy and Uncertainty During Generation}

A growing body of work treats per-token distributional signals (entropy, mutual information, confidence, uncertainty) as zero-cost telemetry for LLM generation quality. These methods all share the observation that the next-token distribution carries information about whether the model is on a healthy trajectory, and they differ chiefly in what signal they extract, what they predict, and how they intervene.

ERGO operationalizes Shannon entropy of the per-token distribution as a drift detector for multi-turn conversations: it computes a delta-entropy ΔH between turns and triggers context resets when ΔH crosses a calibrated threshold τ, yielding a 56.6\% average performance gain over the full-context baseline and a 35.3\% reduction in unreliability across five models and tasks, with strong negative correlation (r = -0.81 on Llama 3.1 70B) between ΔH and task performance \cite{ergo}. Crucially, ERGO ablations show that entropy-guided triggering beats both random and fixed-interval resets, validating entropy as a degradation signal rather than mere noise \cite{ergo}. Semantic Entropy Probes push this further by training linear classifiers on hidden states at a single token position to approximate semantic entropy without the 5-10x cost of multi-sample clustering, and show that hidden states encode semantic uncertainty even before generation begins, generalizing out-of-distribution by +7.7 to +10.5 AUROC over accuracy probes \cite{semanticentropy}. "Demystifying Reasoning Dynamics" measures mutual information (via HSIC) between token representations and correct answers, identifying "thinking tokens" with sharp MI peaks whose suppression collapses GSM8K from 72.7\% to 51.3\% and bounding expected error by 1 - MI(s;Y) (Theorem 1), establishing MI as a principled per-token quality signal \cite{demystifying}. TokUR decomposes token-level uncertainty into aleatoric and epistemic components via low-rank Bayesian weight perturbation, achieving 80.64\% AUROC for hallucination detection on MATH500 while keeping the base weights frozen \cite{tokur}. FineCE complements these by training a supervised confidence estimator with Backward Confidence Integration to produce fine-grained per-position scores throughout generation, addressing LLM overconfidence with continuous rather than coarse signals \cite{mindgeneration}.

HERALD sits squarely in this family but pivots the formulation in three ways. First, where ERGO uses a univariate scalar threshold on entropy and SEPs/TokUR/FineCE produce calibrated uncertainty or confidence scores, HERALD trains a multivariate XGBoost classifier over a vector of zero-cost logit features (entropy, variance, skewness, concentration metrics) treated as a time series. Second, where these methods score the current token, HERALD predicts a binary catastrophe label H tokens ahead, making it a forecaster of compression-induced looping and non-termination rather than a post-hoc uncertainty estimator \cite{tokur,semanticentropy}. Third, where ERGO and SEPs target conversational drift and hallucination under standard inference, HERALD targets the specific failure manifold induced by KV-cache compression. Adjacent but categorically different is UNComp, which also uses entropy at runtime but on hidden states (truncated matrix entropy) to decide which heads and tokens to evict during compression rather than to predict downstream failure: entropy as an eviction policy, not a hazard forecast \cite{uncomp}. HERALD inherits the zero-cost telemetry premise that ERGO, SEPs, and the MI literature legitimize, and repurposes it as the substrate for closed-loop control of compression aggressiveness.

\section{Detecting Failures From Token-Level Signals}

The closest known analogue to HERALD is HALT \cite{halt}, which treats per-token log-probability sequences as time series and feeds them to a lightweight bidirectional GRU (with entropy summary statistics, rank proxies, and temporal differences) to score hallucinations. HALT operates without access to hidden states or external validators, achieving a 30x parameter reduction over fine-tuned ModernBERT baselines and a 60x speed gain over the Lettuce span detector while remaining compatible with black-box proprietary APIs \cite{halt}. The methodological kinship with HERALD is direct: both extract per-token logit signals, both train compact predictors over those signals, and both deliberately avoid hidden-state access so the detector can be deployed alongside any inference stack.

HERALD diverges on three axes. First, the failure mode is qualitatively different. HALT targets factual errors, logical inconsistencies, and reasoning failures in standard inference, with HALT-L reaching 63.04 macro-F1 averaged across HUB capability clusters and showing markedly stronger results on knowledge-heavy clusters than on reasoning-heavy ones, reflecting how unevenly calibration cues track ordinary hallucinations \cite{halt}. HERALD targets compression-induced catastrophic failures (looping, non-termination, instruction amnesia) that arise specifically when KV-cache eviction discards context the decoder still needs; these are not hallucinations of fact but breakdowns of the generation process itself. Second, HERALD uses XGBoost over engineered logit features rather than a GRU over raw log-probability sequences, trading sequence modelling capacity for interpretable feature importances and cheaper inference. Third, HERALD is built around a configurable hazard horizon: it predicts H tokens before catastrophe onset, where HALT scores hallucinations after the fact. This pre-onset framing is what makes HERALD usable as an intervention substrate rather than a post-hoc audit. The cross-model transfer collapse HALT reports (AUROC dropping from 70.02 to 55.24 when the LLaMA-trained variant is applied to Qwen) also flags a constraint HERALD must navigate, since calibration dynamics may similarly fail to transfer across compressors and base models \cite{halt}.

Adjacent work confirms that other token-level signals can flag broken generation, but each differs from HERALD in either signal or objective. Lookback Lens detects contextual hallucinations from attention-map ratios between context and newly generated tokens, training a logistic regression on these features and transferring from a 7B detector to a 13B generator without retraining \cite{lookbacklens}; HERALD instead uses logit distributions and predicts compression-induced catastrophes rather than contextual contradictions. Earlier work on neural text degeneration shows that maximisation-based decoding falls into degenerate repetition through a positive feedback loop in which repeated phrases gain probability with each repetition \cite{curiouscase}, and that threshold-based truncation methods can be characterised through the softmax bottleneck rather than ad hoc cutoffs \cite{closingcurious}. These results establish that per-token distributional properties carry diagnostic signal about broken generation, but they propose decoding strategies, not predictive hazard detectors over a configurable horizon, and they say nothing about KV-cache compression as the failure mechanism HERALD is built to forecast.

\section{Catastrophic Failure Modes in Long-Context and Reasoning LLMs}

HERALD's central premise is that the failures it forecasts (looping, non-termination, instruction drift) are not artifacts of compression alone but well-attested structural failure modes of long-context, autoregressive generation. The literature documents these phenomena along three axes: surface-level degeneration of the output stream, geometric collapse of the reasoning trajectory, and drift of the planning or instruction-following signal across long horizons.

Looping and non-termination are the most overt of these failures. LoopLLM demonstrates that autoregressive LLMs are structurally vulnerable to repetitive generation, with adversarially induced runs reaching over 90\% of the maximum output length compared to roughly 20\% under benign decoding across 12 open-source models \cite{loopllm}. The same work shows that low-entropy, plausible-looking repetition evades simple entropy-based filtering, indicating that surface fluency is preserved even as the generation collapses into a degenerate cycle \cite{loopllm}. This establishes looping and length-saturation as bona fide failure modes in modern decoders, not edge-case artifacts: if adversarial prompts can drive a model into this regime, sufficiently aggressive perturbations of the inference path (such as KV-cache compression) can plausibly do the same without an attacker.

A second strand of work shows that reasoning trajectories themselves carry structure that can degrade. Spectral analysis of attention graphs distinguishes valid from invalid mathematical reasoning with effect sizes |d| >= 2.09 across seven transformer families, reaching 85-95.6\% classification accuracy and up to 95.6\% with calibrated thresholds \cite{geometryreason}. The fact that validity has a measurable internal signature implies, conversely, that invalidity is a structured regime that internal signals can flag, consistent with HERALD's premise that catastrophes leave per-token traces.

A third strand localises where these trajectories break. Thought-anchor analysis shows that a small subset of sentences, typically those involving planning or uncertainty management, exerts outsized counterfactual influence on the rest of a chain of thought, with specialised attention heads consistently routing focus toward them \cite{thoughtanchors}. The preplan-and-anchor work refines this picture: LLMs exhibit a recurring rhythm in which preplan tokens carry +51.97\% higher entropy and anchor tokens receive concentrated future attention influence, and disrupting this rhythm visibly degrades reasoning \cite{preplananchor}. Plan or instruction drift is therefore a documented consequence of perturbing these structural tokens, not a hypothetical mode.

HERALD predicts a subset of these phenomena (specifically those triggered by KV-cache compression) ahead of onset, complementing this descriptive literature with a forecasting capability.

\section{Adaptive and Controller-Based Inference}

A growing line of work treats inference itself as a controllable process: rather than fixing a compression policy ahead of time, these systems modulate eviction, head selection, or layer-wise budgets while the model decodes. The closest precedent to HERALD's intended deployment is ASR-KF-EGR \cite{asrkfegr}, a feedback-loop scheme that does not permanently evict tokens but instead temporarily freezes low-importance KV pairs to off-GPU storage according to a sublinear schedule, with entropy-guided recovery that unfreezes them when context demands shift. ASR-KF-EGR achieves 55-67 percent compression on LLaMA-3 8B while passing needle-in-a-haystack retrieval, demonstrating that reversible, signal-driven KV management can preserve quality. Crucially, the control signal it consumes is an attention-interaction entropy score computed inside the compressor itself; the loop is self-contained and assumes that the chosen attention statistic is a faithful proxy for downstream generation health. HERALD is naturally complementary: a logit-level hazard predictor that can supply the external "is the generation going off the rails" signal that a freeze-and-recover scheme like ASR-KF-EGR otherwise has to infer indirectly.

Several systems generalize this idea to dynamic per-layer budgets that condition on input characteristics. DynamicKV \cite{dynamickv} establishes global and per-layer budgets that update as inference proceeds, retaining as little as 1.7 percent of the cache while preserving roughly 85 percent of full-KV LongBench performance, and surpassing prior methods by 11 points under extreme 0.9 percent compression on Needle-in-a-Haystack. Its policy is driven by task-specific activation patterns observed across layers. CAKE \cite{cake} computes per-layer preference scores from spatial attention dispersion and temporal attention shift, allocating budget proportionally and using an attention-shift-tolerant eviction indicator; it sustains quality at 3.2 percent of the cache and yields over 10x decoding speedup at 128K context. ThinKV \cite{thinkv} specializes the controller to reasoning models, using attention sparsity to identify "thought types" within chain-of-thought and applying hybrid quantization-eviction so that critical reasoning steps stay at high precision while less critical ones are progressively evicted, reaching near-lossless accuracy below 5 percent of cache size with up to 5.8x throughput.

Other controllers operate at the head or block level rather than the layer level. RLKV \cite{whichheads} uses Group Relative Policy Optimization with mixed-attention gating adapters to learn which heads are reasoning-critical, then deploys the resulting policy for selective head-wise eviction, achieving 20-50 percent reduction with negligible loss; notably, its analysis explicitly catalogs three compression-induced failure modes (repetitive, incorrect, overlength errors) that overlap directly with HERALD's looping and non-termination targets, but the failure taxonomy is post-hoc rather than online. Runtime-adaptive pruning \cite{runtimepruning} frames the problem as an MDP over MHA and FFN blocks plus KV-cache, with an RL agent reading runtime state (request mix, memory budget) and a Greedy Sequential Importance scorer; under an 80 percent memory budget on Llama-2-7B it reaches 52.10 perplexity versus 58.67 and 92.41 for static baselines. Lethe \cite{lethe} couples layer-wise sparsity allocation (Hoyer metric) with temporal decay-weighted retention thresholds, achieving 91.7 percent memory reduction and up to 2.56x throughput on reasoning benchmarks.

A different class of controller compresses the generation itself. LightThinker \cite{lightthinker} trains the model to encode verbose chain-of-thought into compact gist tokens via thought-based attention masks, cutting peak tokens by 70 percent on Qwen with only a 1 percent accuracy drop and quantifying the effect with a Dependency metric on historical-token reliance. Adaptive sparse-attention scaffolds \cite{seerattention} and inference-time hyper-scaling policies \cite{hyperscaling} extend the same controller philosophy to attention computation and compute-allocation respectively.

Across all of these systems, the policy is well specified but the runtime quality signal is not: each chooses an internal proxy (attention entropy, sparsity, redundancy, RL-learned head importance, perplexity-on-removal) and trusts it to reflect downstream catastrophic risk. None of them measures whether generation is actually about to collapse. HERALD targets exactly this gap. It does not propose another control law; it treats the compressor as a black box and supplies a calibrated, horizon-aware hazard probability derived from per-token logit dynamics. The natural composition is straightforward: a controller from this family selects the action (loosen, tighten, freeze, unfreeze, reallocate), and HERALD provides the closed-loop measurement that tells it when to act before failure manifests, not after.

\section{Empirical Limits and Failures of KV-Cache Compression}

A growing body of empirical work establishes that KV-cache compression does not degrade gracefully. Across compressors, ratios, tasks, and model families, the literature converges on a picture in which compression triggers sharp, structured failures that current pipelines neither anticipate nor monitor. We organize this evidence by failure mechanism.

\textbf{Eviction fragility and mean-aggregation traps.} \cite{defensivekv} documents that the standard mean-pooled importance score used by attention-based eviction policies is unexpectedly fragile: a single outlier token can devastate downstream quality even when the average importance is preserved, and their defensive aggregation strategy reduces quality loss by 2.3x to 4.3x at 20\% cache size \cite{defensivekv}. The fragility is not aesthetic; \cite{foresightkv} quantifies the asymmetric cost of evicting the wrong tokens, showing that evicting low-entropy tokens (precisely the tokens that look "safe" by entropy heuristics) causes loss increases of +147\% in Math, +75\% in Code, and +187\% in Summarization relative to random eviction \cite{foresightkv}. The implication is that the importance signals current evictors rely on are systematically miscalibrated against long-horizon use. \cite{holdonto} sharpens this by showing that even among heavy-hitter trackers, eviction strategy interacts non-trivially with reasoning: lower cache budgets trigger longer reasoning traces, exposing a hidden tradeoff between memory and inference cost in which the lowest-performing strategy (KNorm) causes the greatest output elongation, an emergent failure mode invisible to any single-pass accuracy metric \cite{holdonto}.

\textbf{Compression breaks reasoning and arithmetic.} \cite{llmabilities} establishes that capability degradation under compression is task-dependent and discontinuous rather than smooth. On long-context generation in arithmetic reasoning, performance declines by more than 20\% at compression ratios below 30\%, with precision dropping from approximately 0.75 to below 0.5 once the budget crosses a 20\% threshold \cite{llmabilities}. Shorter prompts are more vulnerable still: in 1-shot settings, performance collapses from 0.5 to 0.05 between 30\% and 10\% compression, while 8-shot settings degrade more gracefully \cite{llmabilities}. The pattern, capabilities surviving down to a knee then falling off a cliff, is characteristic of a hazard process rather than a Gaussian noise model. \cite{quantreasonhurts} reports the same shape for the adjacent compression family of quantization: W8A8 and W4A16 are effectively lossless for reasoning models, but lower bit-widths introduce significant accuracy risks, with model size, origin, and task difficulty as decisive variables \cite{quantreasonhurts}. Compression breaks reasoning when it breaks, and the break is sudden.

\textbf{Learned importance does not beat simple heuristics.} A natural response is to train a stronger token-importance predictor. \cite{limitslearned} shows this hits an information-theoretic ceiling. A 1.7M-parameter Speculative Importance Predictor shows no statistically significant advantage over random selection across datasets and ratios, and is in fact dominated by position-based heuristics at 10-25\% retention and by prefill attention at 50-75\% retention \cite{limitslearned}. The mutual-information analysis explains why: position alone provides 0.31 +/- 0.05 bits about future token importance, prefill attention 0.28 +/- 0.04 bits, while key vectors (the substrate learned scorers consume) provide only 0.12 +/- 0.03 bits \cite{limitslearned}. The signal needed to evict optimally is largely absent from the features available pre-generation, so building better evictors is a path with diminishing returns.

\textbf{Systemic pitfalls and benchmark-level evidence.} \cite{pitfalls} catalogs six distinct failure modes including instruction degradation, eviction-induced instruction ignoring, and system prompt leakage; ROUGE-L overlap between leaked output and protected system prompts rises sharply under StreamingLLM compression, indicating progressive disclosure of instructions as the ratio increases \cite{pitfalls}. At the benchmark level, \cite{scbench} reports that at a 1/32 compression ratio StreamingLLM achieves only 15.5\% average performance versus a 48.7\% baseline on Llama-3.1-8B, and that sub-O(n) memory methods perform well on first queries but lose accuracy on subsequent multi-turn requests, exposing failure modes invisible to single-shot evaluation \cite{scbench}. \cite{whatmustgive} complements this with a 10-method, 16-task sweep concluding that prefill compression is critical but risky, and that token-dropping methods sometimes excel and sometimes severely underperform on the same task family, while needle-in-a-haystack retrieval remains challenging for all methods \cite{whatmustgive}. Reliability profiles differ fundamentally across method families and resist any single fixed configuration.

Taken together, these results establish that compression failures are real, frequent, mechanism-diverse, and not predictable from benchmark-averaged accuracy curves: a model can be operating in a "lossless" regime by aggregate metrics while individual generations loop, leak, or stall. Yet none of this work provides a runtime hazard predictor over a generation in progress; the literature characterizes where compression breaks and proposes static method-level mitigations, but leaves the online-monitoring problem open. HERALD fills that gap by treating any compressor as a black box and forecasting catastrophe onset from per-token logit features.
